\section{片上物化重放与上下文调度}
\label{sec:materialize}

\subsection{一次物化与多轮重放}

一个空间 tile 到达 LASA 后，窗口生成单元按 HWC 顺序消费输入，同时把全部
$P_l$ 轮的 18 元素向量写入缓存。第一组向量可在后续行仍在生成时开始计算；
当相应地址发布后，当前输出通道块依次读取 $P_l$ 轮，随后下一输出通道块从同一
缓存起点重放。输入 tile 的生命周期因此覆盖 $Q_l$ 个输出通道块，而 DDR 只
参与生命周期的起点。

这一机制把数据复用分成两个层次：窗口生成利用输入行和相邻输出位置的空间
重用，重放缓存利用不同输出通道权重共享同一窗口。前者避免软件 im2col，后者
实现式\eqref{eq:dlasa} 的片外流量。缓存不保存完整网络张量，只保存当前和下一
空间 tile 的阵列向量，其容量随 $P_l$ 和 tile 像素数变化。

\subsection{层自适应空间分块}

设候选 tile 高度为 $h$，输出宽度为 $W_o$。分块同时受部分和像素、物化向量、
HWC 打包和输入行宽约束：
\begin{align}
W_oh &\le 1024,\label{eq:tile-psum}\\
W_ohP_l &\le 32768,\label{eq:tile-material}\\
W_oh\left\lceil C_{out}/32\right\rceil &\le4096,\label{eq:tile-pack}\\
\left\lceil C_{in}/4\right\rceil\left\lceil W_i/2\right\rceil&\le2048.
\label{eq:tile-line}
\end{align}
第一式限制部分和像素，第二式限制物化缓存条目，第三式限制紧凑写回，第四式
限制输入行存储。满足容量后，选择能够减少 tile 数并控制参数重复的最大安全
高度。最终 13 层高度为 2、4、8、8、13、13、8、13、13、13、6、13、13，
对应\cref{tab:layer-config}。

m13 的 $P_l=256$，使用 8 行 tile 后物化条目为 $13\times8\times256=26624$；
m19 的 $P_l=192$，6 行 tile 已达到 $26\times6\times192=29952$，若增至 7 行
将超过 32768 条容量。因此层自适应并非改变阵列规模，而是在同一物化重放
机制下选择可容纳且启动次数较少的空间粒度。早期宽特征图使用小 tile，深层
则尽量覆盖完整高度。

\subsection{阵列参数与分块的协同}

设阵列行数为 $R$，双输出处理单元列数为 $C$，则并行乘加数、计算轮次和输出
通道块分别为
\begin{equation}
\begin{aligned}
N_{MAC}&=2RC,\\
P_l&=\left\lceil\frac{K_l}{R}\right\rceil,\qquad
Q_l=\left\lceil\frac{C_{out,l}}{2C}\right\rceil .
\end{aligned}
\label{eq:array-space}
\end{equation}
增加 $R$ 可以减少长层计算轮次，却会线性增加 DSP、输入广播宽度和每个权重
上下文；增加 $C$ 可以减少输出块数，却会扩大部分和 packet、重量化 lane 和
写回网络。空间 tile 高度 $h$ 不改变阵列峰值，但同时影响缓存深度、dispatch
次数和行边界重复。

\begin{table}[t]
\centering
\caption{三个主要设计参数的收益与代价。}
\label{tab:design-space}
\small
\setlength{\tabcolsep}{3pt}
\begin{tabular}{lp{25mm}p{34mm}}
\toprule
参数 & 直接收益 & 主要代价 \\
\midrule
行数 $R$ & 减少 $P_l$ & DSP和广播宽度增加，权重上下文变宽 \\
列数 $C$ & 减少 $Q_l$ & 部分和、重量化和写回网络变宽 \\
tile高度 $h$ & 减少dispatch & 窗口、部分和及打包容量增加 \\
\bottomrule
\end{tabular}
\end{table}

本文选择 $R=18,C=16$，形成 576-MAC 阵列和 32 通道输出块。该规模能够在
XCK26 上为窗口缓存保留 48 个 URAM，并使两个 255 通道检测头只产生一个无效
lane。对于 m13，继续增加行数虽能减少 256 个轮次，却会同步扩大 32 个输出块
的权重存储和跨列广播；对于早期层，增加行数对原本很短的规约维收益有限。
因此，层自适应主要通过 $h$ 调节空间粒度，$R$ 和 $C$ 保持为所有层共享的物理
结构。该分工把运行时变化限制在描述参数和地址范围内，有利于保持统一时序。

从容量式\eqref{eq:tile-psum}--\eqref{eq:tile-line}可见，$R$ 还通过 $P_l$
间接影响物化条目。理想情况下 $RP_l$ 接近 $K_l$，总向量字节随 $R$ 变化较小；
实际设计中每一轮都有有效位、地址和提交开销，较大的 $R$ 会降低轮次控制开销，
却提高单轮扇出和尾轮浪费。18 行在完整网络中兼顾 m0 的 27 元素规约和 m13 的
4608 元素规约，使短层保持少量轮次、长层仍能由片上重放覆盖。

\subsection{上下文调度}

计算上下文由层、空间 tile、输出通道块和 $K$ 维轮次确定。LASA 为权重和输入
缓存维护当前与下一上下文。当当前阵列计算时，空闲权重缓冲接收下一轮参数；
窗口生成单元可在活跃读者不再访问相关行后开始生成下一 tile。控制器只在输入
向量、权重和部分和三者均就绪时启动计算，从而使式\eqref{eq:tlasa} 中的阶段
重叠具有明确的资源条件。

部分和缓冲采用双槽反馈。非末轮结果写入返回槽，下一轮计算从读槽获得同一
像素和输出通道块的部分和；读写指针在完整 packet 提交后交换。该结构允许
结果排出与下一组输入准备并行，并把宽部分和数据与下游就绪控制隔离。最后一轮
结果进入重量化后释放部分和槽，下一上下文即可取得该资源。

当当前上下文的最后一个结果、下一权重和下一输入均已就绪时，控制器直接切换
上下文，无需经过全局空闲状态。该切换只更新窄控制状态，数据缓冲的释放仍以
提交事件为准。因此，权重提前装载缩短 $T_W$，双槽反馈缩短 $T_F$，上下文直接
切换缩短 $T_{tail}$；三者都服务于片上物化重放形成的连续输入序列。

\begin{table}[t]
\centering
\caption{上下文调度机制、隐藏目标及启动条件。}
\label{tab:schedule-mechanisms}
\small
\setlength{\tabcolsep}{2pt}
\begin{tabular}{p{18mm}p{13mm}p{38mm}}
\toprule
机制 & 目标阶段 & 必要条件 \\
\midrule
窗口提前生成 & $T_M$ & 待写行不与活跃读者重叠 \\
权重提前装载 & $T_W$ & 下一权重槽空闲且身份匹配 \\
双槽反馈 & $T_F$ & 返回 packet 完整提交 \\
上下文直接切换 & $T_{tail}$ & 下一输入、权重和资源同时就绪 \\
\bottomrule
\end{tabular}
\end{table}

\cref{tab:schedule-mechanisms} 将重叠机会与资源条件对应起来。调度器不会根据
预测时间提前覆盖数据，而是等待能够在硬件中观察的提交事件。这样，性能重叠
由缓冲容量和依赖关系自然产生：长层拥有更多稳态轮次，权重装载和反馈更容易
被计算覆盖；短层的稳态区较短，直接切换和减少 tile 数更重要。层自适应分块
与上下文调度由此形成互补关系。

\begin{figure*}[t]
\centering
\placeholderbox{52mm}{视觉占位：串行时间线依次显示输入、权重、计算、反馈；LASA 时间线显示一次物化、权重提前装载、连续部分和反馈和上下文直接切换}
\caption{串行逐轮执行与 LASA 上下文调度的时间线。片上重放消除重复输入阶段，
独立缓冲使权重、计算和部分和反馈在稳态中重叠。}
\label{fig:timeline-placeholder}
\end{figure*}

\subsection{正确性不变量}

重叠调度遵循三条外部可验证的不变量。

\begin{enumerate}[leftmargin=2.1em]
  \item \textbf{缓冲区唯一占用：}任一权重、窗口或部分和槽在同一周期只属于
  一个计算上下文；下一上下文只有在释放事件到达后取得该槽。
  \item \textbf{数据不覆盖：}窗口写地址不能越过活跃读者仍需访问的最小地址；
  遇到反压时写入延后，已发布的向量保持稳定。
  \item \textbf{顺序提交：}输出顺序由空间 tile、输出通道块、像素和通道的
  固定次序决定；内部预载和完成顺序不能改变软件可见 HWC 字节序列。
\end{enumerate}

三条不变量分别约束空间、时间和外部次序。它们使重叠机制可以通过随机反压、
尾轮、尾通道和完整网络结果进行检验，而无需观察内部状态机。本文的逐节点
验证即以紧凑 HWC 输出为边界，将每个硬件节点与整数语义参考逐字节比较。

唯一占用与不覆盖还给出一个直接推论：任一已发布窗口向量在最后一个读者释放
之前保持不变。因此，不同输出通道块读取同一地址时得到相同输入字节，片上
重放与从不可变输入张量重复读取具有相同语义。顺序提交则给出第二个推论：
即使下一上下文先完成权重或输入准备，外部输出仍按 tile、输出块、像素和通道
的字典序形成。两条推论把内部并行性约束为确定的外部 HWC 流，使主机参考无需
模拟具体周期即可验证结果。

反压场景下，输入、反馈和输出三条链路可能在不同周期停止。LASA 将有效数据与
占用状态同时保持，恢复后从同一提交点继续；异常状态则阻止新上下文取得资源。
因此，随机停顿主要改变完成时间，不改变上下文身份、窗口内容和输出顺序。这一
性质是从单层测试扩展到完整 DAG 和长时间运行的基础。

\subsection{吞吐与利用率定义}

设完整网络运算量为 $M$ MAC，PL 执行时间为 $T_{PL}$，阵列理论峰值为
$576f$ MAC/s。按一次乘法和一次加法计为两个操作，定义
\begin{equation}
P_{eff}=\frac{2M}{T_{PL}},\qquad
U_{array}=\frac{P_{eff}}{2\times576f}.
\label{eq:utilization}
\end{equation}
$P_{eff}$ 包含尾轮、尾通道、空间 tile 边界和调度等待，是从完整网络观测的
有效吞吐；$U_{array}$ 将其与固定阵列峰值比较。硬件同时记录有效计算周期、
输入等待、权重等待、部分和反馈间隔和结果排出等待，用于区分计算量差异和
数据供应空泡。
